\begin{frame}{Strategia deterministyczna}
    Gracze mogą ustalić strategię. Ich strategia to para funkcji 
    \[f,g:\set{0,1} \to \set{0,1}\]
    Alicja używa funkcji \(f\), Bartosz funkcji \(g\) zatem:
    \[f(x)=a \qquad \quad g(y)=b\]
    Warunek końcowy to:
    \[x\cdot y=f(x)+g(y) \mod{2}\]
\end{frame}
\begin{frame}{Optymalna strategia deterministyczna}
    Warunek końcowy to:
    \[x\cdot y=f(x)+g(y) \mod{2}\]
    Mamy 4 możliwe wejścia
\only<1>{
\begin{align*}
    (0,0): \qquad f(0)+g(0)=0\\
    (0,1): \qquad f(0)+g(1)=0\\
    (1,0): \qquad f(1)+g(0)=0\\
    (1,1): \qquad f(1)+g(1)=1
\end{align*}
}
\only<2>{
\[
\begin{array}{rcl}
(0,0): & f(0)+g(0)=0 & \\
(0,1): & f(0)+g(1)=0 & \\
(1,0): & f(1)+g(0)=0 & \\
(1,1): & f(1)+g(1)=1 & \\
\end{array}
\;\left.\vphantom{\begin{array}{c}
(0,0)\\(0,1)\\(1,0)\\(1,1)
\end{array}}\right|\;
\begin{array}{c}
\\
\displaystyle +\\
\\
\end{array}
\]
}
\onslide<2>{
\[2[f(0)+f(1)+g(0)+g(1)]=1\]
}
\end{frame}
